% ======================================================================
% Chapter 1 -- From Building Agents to Designing Solutions
% Designing Agentic AI Solutions, Chris Kuo
% Part A: opening, the design gap, the autonomy spectrum
% ======================================================================

\chapter{From Building Agents to Designing Solutions}
\label{ch:foundations}

\begin{objectives}
After reading this chapter, you will be able to:
\begin{enumerate}[leftmargin=*,itemsep=2pt]
  \item Explain the difference between \emph{building} an agent and \emph{designing} a solution, and name the five design skills this book trains.
  \item Place a proposed system on the autonomy spectrum, from fixed rules to multi-agent systems, and argue for the least autonomy that meets the business goal.
  \item Describe the four building blocks of an agentic solution---plain code, retrieval-augmented generation (RAG), the agent, and the Model Context Protocol (MCP)---and say what each one contributes and how each one fails.
  \item Decompose a business process into steps and assign each step to a building block.
  \item Set up and run a minimal agent with Google's Agent Development Kit (ADK) that calls a tool served by an MCP server.
  \item Name the Google Cloud services that support each stage of the production lifecycle used throughout the book.
\end{enumerate}
\end{objectives}

% ----------------------------------------------------------------------
\section{Two Demos on a Friday Afternoon}
\label{sec:two-demos}

Harbor Industrial Supply is a regional distributor of pumps, valves and fittings.\footnote{Harbor Industrial Supply and all figures attached to it are invented for teaching. The pattern is common across business-to-business distributors.} Its sales representatives cannot quote large orders on their own. Any order above \$25{,}000, or any request for a discount off the list price, goes to a six-person pricing desk. The desk checks the customer's contract terms, current stock, the cost of each item and the company's discount policy, then sends back a quote. About 400 requests reach the desk each week. The median request waits 26 hours for an answer, and the vice president of sales believes the company loses deals to faster competitors because of it.

The vice president asked two internal teams to show what AI could do. Each team had two weeks.

\paragraph{Team A} built an agent. They gave a large language model a long system prompt that described the discount policy, and connected it to fourteen tools: product search, price lookup, inventory, customer lookup, contract lookup, margin calculator, email sender, CRM updater and six more. In the demo, a sales representative typed, ``Quote 40 units of the 3-inch stainless ball valve for Delta Water, they want 18\% off.'' Within a minute the agent looked up the customer, found the contract, checked stock, computed a price and drafted an email. The room applauded.

Then the finance director asked three questions. What stops the agent from approving 18\% when policy caps it at 15\% without a manager's sign-off? Where does the policy text come from, and what happens when it changes next quarter? How much does each quote cost to produce, and how will we know if the quotes are wrong? Team A had partial answers. The cap was ``in the prompt.'' The policy was ``in the prompt.'' The cost was ``probably a few cents.'' Nobody had measured the error rate, because there was nothing to compare against.

\paragraph{Team B} showed a working system too, but it was smaller. Before the demo, they put up a single diagram (Figure~\ref{fig:quote-sketch} shows a version of it). It showed which steps were handled by plain code (the discount cap, the margin floor, the approval threshold), which used retrieval over the policy documents (with the policy version cited on every quote), which used a language model (reading a free-text request and writing the justification a manager would see), and which called systems of record through a small set of tools. A manager approved every discount above 15\%. Beside the diagram was a table comparing their design with a single all-purpose agent like Team A's, on cost, speed, error risk and auditability. Below it was one page addressed to the vice president: the target turnaround (under four hours), the cost per quote, the expected effect on win rate, the main risks, and a four-week pilot with go/no-go criteria.

\begin{figure}[htbp]
\centering
\begin{tikzpicture}[
  node distance=7mm and 12mm,
  box/.style={draw, rounded corners=2pt, align=center, font=\footnotesize, minimum height=9mm, text width=2.7cm},
  agent/.style={box, fill=agentfill},
  code/.style={box, fill=codefill},
  mcp/.style={box, fill=mcpfill},
  rag/.style={box, fill=ragfill, shape=cylinder, shape border rotate=90, aspect=0.18, rounded corners=0pt},
  human/.style={box, fill=humanfill, shape=rectangle, double},
  sor/.style={box, fill=white, dashed},
  arr/.style={-{Stealth[length=2mm]}, thick},
  lab/.style={font=\scriptsize, fill=white, inner sep=1pt}
]
% main flow
\node[box, fill=white] (rep) {Sales rep\\(email / chat)};
\node[agent, right=of rep] (extract) {Extract request\\(agent, Level 1)};
\node[code, right=of extract] (validate) {Validate\\schema (code)};
\node[mcp, below=14mm of validate] (erp) {ERP MCP server\\{\scriptsize\ttfamily get\_item\_pricing}\\{\scriptsize\ttfamily get\_contract\_terms}};
\node[rag, left=of erp] (kb) {Policy \&\\contract text\\(RAG)};
\node[code, below=14mm of erp] (policy) {Apply policy:\\cap, margin floor\\(code)};
\node[human, left=of policy] (mgr) {Manager\\approval};
\node[agent, below=12mm of policy] (draft) {Draft cover note\\(agent); numbers\\filled by code};
\node[mcp, left=of draft] (crm) {CRM MCP server\\{\scriptsize\ttfamily create\_quote\_draft}};
\node[human, left=of crm] (send) {Rep reviews\\and sends};
\node[sor, right=of erp] (erpdb) {ERP\\(system of record)};
\node[sor, right=of draft] (crmdb) {CRM\\(system of record)};

\draw[arr] (rep) -- (extract);
\draw[arr] (extract) -- (validate);
\draw[arr] (validate) -- node[lab, right] {complete} (erp);
\draw[arr] (validate.north) to[bend right=35] node[lab, above] {missing info: ask rep} (rep.north);
\draw[arr] (erp) -- (policy);
\draw[arr] (kb) -- node[lab, above, sloped, pos=0.4] {policy version} (policy);
\draw[arr] (policy) -- node[lab, above] {exception} (mgr);
\draw[arr] (mgr.south east) -- node[lab, sloped, above, pos=0.5] {approved} (draft.north west);
\draw[arr] (policy) -- node[lab, right] {within policy} (draft);
\draw[arr] (draft) -- (crm);
\draw[arr] (crm) -- (send);
\draw[arr, dashed] (erp) -- (erpdb);
\draw[arr, dashed] (crm.south) to[bend right=20] (crmdb.south);
\end{tikzpicture}
\caption{A first, informal sketch of Team B's quote design. Shading marks the block type: agent steps, code, MCP servers, retrieval, and double borders for people. Chapter~\ref{ch:toolkit} introduces the standard notation.}
\label{fig:quote-sketch}
\end{figure}

The vice president funded Team B.

Team A's engineers were at least as skilled as Team B's. Their agent was impressive, and with a few more weeks it might have worked well. The difference was that Team B had \emph{designed a solution}: they decided what the system should and should not do, chose a building block for each part on purpose, compared alternatives, and tied the whole thing to a business result someone could check. Team A had \emph{built an agent} and hoped the business case would follow.

This book is about doing what Team B did, repeatably, for real business problems.

% ----------------------------------------------------------------------
\section{The Design Gap}
\label{sec:design-gap}

\subsection{Building is getting cheap; designing is not}

A few years ago, getting a language model to call a function, retrieve from a document store or loop until a task was done took real engineering. Today a student can do all three in an afternoon. Open-source frameworks such as Google's Agent Development Kit, LangGraph, the OpenAI Agents SDK and the Claude Agent SDK hide the loop. Vendors publish sample agents for customer service, data analysis and document processing that deploy with one command. Coding assistants write much of the remaining code.

When building gets cheaper, the bottleneck moves. The hard questions are no longer ``Can we make it call the inventory API?'' but:
\begin{itemize}[itemsep=1pt]
  \item Which parts of this process should a model decide at all?
  \item What must never be left to a model, and how do we enforce that?
  \item Where does the system's knowledge come from, who is allowed to see it, and how does it stay current?
  \item Which actions can the system take on its own, and which need a person?
  \item What does it cost per transaction, what does a mistake cost, and does it pay off?
  \item How will we know, in production, whether it is working?
\end{itemize}
These are design questions. Courses and tutorials rarely teach them, because they tend to cover RAG as a pipeline, agents as a loop and MCP as a protocol, each on its own. Vendor samples rarely answer them either. A sample shows \emph{one} working design; it seldom says what the alternatives were, why this one was chosen, or whether it would pay off for your business. That is not a criticism of the samples. Their job is to show how the framework works. Your job, as a designer, is different.

\subsection{Five design skills}

We will call the distance between ``can build an agent'' and ``can design a solution'' the \keyterm{design gap}. Closing it takes five skills. Every case in this book exercises all five.

\begin{description}[leftmargin=0pt,itemsep=3pt]
  \item[Decomposition.] Breaking a business process into steps and deciding, for each step, whether it needs \emph{judgment} (a model), \emph{knowledge} (retrieval), \emph{action} (a tool that touches a system), or a \emph{fixed rule} (plain code). Most bad designs fail here first: they hand a model a decision that a three-line \texttt{if} statement would make better, or they hard-code a judgment that really needs a model.
  \item[Boundaries.] Deciding where one component ends and another begins: what becomes an MCP server, which capabilities are tools, which are read-only resources, and which are prompt templates; which agent owns which decision; where a person must step in.
  \item[Trade-offs.] Comparing two or more designs on fixed criteria---quality, cost, speed, reliability, security and privacy, control and auditability, maintainability, and human workload---and writing down the decision and what would change it.
  \item[Production.] Carrying a design through evaluation, deployment, security, observability and operation after launch. A design that cannot be evaluated or monitored is not finished.
  \item[Business justification.] Tying the design to key performance indicators (KPIs), unit economics, the cost of errors, risk and a rollout plan, and explaining all of it to a sponsor in one page.
\end{description}

\subsection{What a design is}

Throughout the book we use the word \emph{design} in a specific sense.

\begin{definition}
A \keyterm{solution design} is a set of decisions about how a business process will be carried out by people, software and AI components, together with the reasons for those decisions, the alternatives that were rejected, and the measurements that will show whether the decisions were right.
\end{definition}

A design is not a diagram, although it usually includes one. It is not code, although code implements it. The definition has three parts on purpose. The \emph{decisions} say what the system does. The \emph{reasons and alternatives} let someone else check the thinking and change it later when conditions change. The \emph{measurements} turn the design into a claim that can be tested in production.

The definition leads directly to the four deliverables you will produce for every case in this book:
\begin{enumerate}[itemsep=1pt]
  \item a \keyterm{solution graph}: a diagram, in a fixed notation, of the components and how work flows between them (Chapter~\ref{ch:toolkit} introduces the notation);
  \item a \keyterm{trade-off table}: at least two alternative designs scored on the criteria above, with a written decision record;
  \item a \emph{working deployment}: the chosen design built with ADK and running on Google Cloud;
  \item a \keyterm{one-page memo} to the executive sponsor: the business case, the risks, and a rollout plan with go/no-go criteria (Chapter~\ref{ch:business} shows how to write it).
\end{enumerate}

Team B, in the opening story, brought all four. That is why they won.

% ----------------------------------------------------------------------
\section{How Much Autonomy? From Rules to Agents}
\label{sec:autonomy}

The word \emph{agent} is used loosely. Marketing material calls nearly any system with a language model in it an agent. To design well, we need sharper terms. The most useful question to ask about any proposed system is: \emph{who decides what happens next?} In a traditional program, the programmer decides, in advance, by writing the control flow. In a fully autonomous agent, the model decides at run time, step by step. Between the two lies a range of options.

\begin{definition}
An \keyterm{agent}, in this book, is a software component in which a language model chooses, at run time, which action to take next---which tool to call, with what arguments, or whether to stop---based on the goal it was given and the results of its earlier actions.
\end{definition}

A system can contain language models without being an agent in this sense. A pipeline that always retrieves five documents and then asks a model to summarize them uses a model, but the programmer fixed the steps. Anthropic's engineering guidance draws the same line between \emph{workflows}, where code orchestrates the model along predefined paths, and \emph{agents}, where the model directs its own process and tool use, and it recommends starting with the simplest option that works \citep{anthropic2024effective}. We adopt the same stance and make it a rule.

Table~\ref{tab:autonomy} lays out six levels. The levels are not a ranking of quality. Each level up buys flexibility and pays for it in cost, latency, unpredictability and difficulty of testing.

\begin{table}[htbp]
\centering
\small
\caption{The autonomy spectrum. Each level up gives the model more control over what happens next.}
\label{tab:autonomy}
\begin{tabularx}{\textwidth}{@{}c>{\raggedright\arraybackslash}p{2.6cm}>{\raggedright\arraybackslash}X>{\raggedright\arraybackslash}X@{}}
\toprule
Level & Name & Who decides the next step & Example at Harbor's pricing desk \\
\midrule
0 & Rules only & Code, fully & Reject any discount above 15\% that lacks a manager's approval. \\
1 & Single model call & Code; the model fills one step & Extract product, quantity and requested discount from a free-text email. \\
2 & Fixed workflow & Code runs a fixed sequence that includes model steps & Extract $\rightarrow$ look up prices $\rightarrow$ apply policy $\rightarrow$ draft the quote text. \\
3 & Routed workflow & A model picks one of several fixed paths & Classify the request as a standard quote, a contract renewal or a complaint, and send it down the matching path. \\
4 & Tool-using agent & The model, in a loop, picks tools until done & Investigate why a customer's contract price disagrees with the list price, calling lookups as needed. \\
5 & Multi-agent system & Several agents, each with its own goal, hand work to one another & A sales-side agent negotiates with a supplier-side agent to secure stock for a rush order. \\
\bottomrule
\end{tabularx}
\end{table}

Two observations follow from the table.

First, \emph{a real solution mixes levels}. Harbor's quote system, as Team B designed it, is mostly Level~2, with a Level~1 extraction step, Level~0 rules for the discount cap and margin floor, and a Level~4 agent only for the minority of requests that need investigation. Asking ``Is this an agent or not?'' about the whole system is less useful than asking it about each step.

Second, \emph{autonomy should be earned}. We will use the following principle in every case.

\begin{principle}[Least sufficient autonomy]
For each step in a process, choose the lowest level of autonomy that meets the step's quality target. Move a step up a level only when you can name the requirement the lower level fails to meet, and when you can evaluate the higher level.
\end{principle}

The principle is not a prejudice against agents. Some steps genuinely need one: the inputs vary too much to enumerate the paths, the number of steps is not known in advance, or the cost of a fixed workflow that handles every case would exceed the cost of an agent that figures it out. Chapter~\ref{ch:business} shows how to put numbers on that comparison. The principle simply says that the burden of proof sits with the higher level, because the higher level is harder to test, harder to audit and usually more expensive to run.
